\documentclass{article}
\usepackage[utf8]{inputenc}
\usepackage[T1]{fontenc}
\usepackage{lmodern}
\usepackage[polish]{babel}
\usepackage{amsmath}
\usepackage{tikz}
\usepackage{algorithm}
\usepackage{algpseudocode}
\usepackage{hyperref}
\usepackage{float}
\usepackage{graphicx}
\usepackage{mathtools}
\usepackage{amsmath}
\usepackage{amsfonts}
\usepackage{amsmath}
\usepackage{amsmath}
\usepackage{booktabs}
\usepackage{subcaption}
\usepackage[margin=1in]{geometry}




\title{Algorytmy Metaheurystyczne - Lista 1}
\author{Wojciech Typer}
\date{}

\begin{document}
\maketitle

\section*{Wprowadzenie - Algorytm Local Search}
Algorytm Local Search (LS) jest jednym z podstawowych algorytmów metaheurystycznych, który służy do rozwiązywania problemów optymalizacyjnych. 
Jego główną ideą jest iteracyjne przeszukiwanie przestrzeni rozwiązań, zaczynając od pewnego początkowego rozwiązania i stopniowo poprawiając 
je poprzez lokalne modyfikacje. Algorytm ten jest prosty w implementacji i często stosowany w praktyce, zwłaszcza w przypadkach, 
gdy przestrzeń rozwiązań jest duża i nie można jej przeszukać w całości.

\section*{Problem komiwojażera (TSP)}
Problem komiwojażera (TSP) jest klasycznym problemem optymalizacyjnym, który polega na znalezieniu najkrótszej możliwej trasy, 
która odwiedza każde z zadanych miast dokładnie raz i wraca do punktu startowego. TSP jest problemem NP-trudnym, co oznacza, 
że nie istnieje znany algorytm, który mógłby rozwiązać go w czasie wielomianowym dla wszystkich przypadków. 
Algorytm Local Search może być używany do znajdowania przybliżonych rozwiązań dla TSP, poprzez iteracyjne poprawianie aktualnego rozwiązania 
poprzez lokalne modyfikacje, takie jak zamiana kolejności odwiedzanych miast.

\newpage

\section*{Zadanie 1}
% WYNIKI do wykresu:
% Oman avg_res: 93848.80 avg_num_steps: 2788.60 best_res: 92363
% Dżibuti avg res:6828.20 avg_num_steps: 34.40 best_res: 6656
% Irlandia avg_res: 229621.5661093335 avg_num_steps: 12099.333333333334 best_res: 228986.43866100002
% Egipt avg_res: 189097.36425066614 avg_num_steps: 10382.333333333334 best_res: 186587.64415899917
% Tanzania avg_res: 197495.40 avg_num_steps: 4232.20 best_res: 196949
% Kanada avg_res: 1416488.80 avg_num_steps: 6797.60 best_res: 1405432
\begin{figure}[h!]
    \centering
    % Pierwszy wykres
    \begin{subfigure}[b]{0.45\textwidth}
        \centering
        \includegraphics[width=\textwidth]{/home/wojteq18/Uni-6/Algorytmy_Metaheurystyczne/Lab/Lista_1/charts/11.png}
        \caption{Średnie rozwiązanie}
        \label{fig:wykres1}
    \end{subfigure}
    \hfill % Rozsuwa wykresy na całą szerokość tekstu
    % Drugi wykres
    \begin{subfigure}[b]{0.45\textwidth}
        \centering
        \includegraphics[width=\textwidth]{/home/wojteq18/Uni-6/Algorytmy_Metaheurystyczne/Lab/Lista_1/charts/12.png}
        \caption{Średnia ilość kroków}
        \label{fig:wykres2}
    \end{subfigure}
    \hfill
    % Trzeci wykres
    \begin{subfigure}[b]{0.45\textwidth}
        \centering
        \includegraphics[width=\textwidth]{/home/wojteq18/Uni-6/Algorytmy_Metaheurystyczne/Lab/Lista_1/charts/13.png}
        \caption{Najlepszy rezultat}
        \label{fig:wykres3}
    \end{subfigure}
    
    \caption{Porównanie wyników algorytmu dla wybranych państw}
    \label{fig:wszystkie_wykresy}
\end{figure}

\subsection*{Analiza wyników}

Analiza uzyskanych wyników pozwala na wyciągnięcie wniosków dotyczących charakterystyki działania algorytmu Local Search w zależności od skali problemu. Zaobserwowane różnice w wynikach są w pełni zgodne z intuicją i wynikają bezpośrednio ze specyfiki badanych instancji:

\begin{itemize}
    \item \textbf{Skala wartości funkcji celu:} Wyższe wartości parametrów \textit{avg\_res} oraz \textit{best\_res} dla państw takich jak Kanada, w porównaniu do znacznie niższych wyników dla Dżibuti, nie świadczą o mniejszej skuteczności algorytmu. Wynikają one bezpośrednio z rozmiaru problemu – większa liczba miast rozmieszczonych na większym obszarze geograficznym naturalnie generuje wyższy koszt całkowity (sumę odległości), co jest zjawiskiem spodziewanym.
    
    \item \textbf{Złożoność obliczeniowa a liczba kroków:} Liczba kroków potrzebna do osiągnięcia poprawy (\textit{avg\_num\_steps}) rośnie wraz ze wzrostem liczby miast w danej instancji. Jest to spowodowane większą złożonością przestrzeni poszukiwań – im więcej miast posiada dany kraj, tym więcej istnieje potencjalnych możliwości zamiany połączeń i ścieżek poprawy, co naturalnie wydłuża proces zbieżności algorytmu do lokalnego optimum.
    
    \item \textbf{Charakterystyka instancji:} Przypadek Dżibuti, charakteryzujący się bardzo małą liczbą kroków i niskim wynikiem, potwierdza, że dla niewielkich zbiorów danych algorytm Local Search błyskawicznie przeszukuje dostępną przestrzeń rozwiązań. Z kolei instancje takie jak Kanada, Egipt czy Irlandia wymagają znacznie większego nakładu pracy, co odzwierciedla ich wysoką złożoność kombinatoryczną.
    
    \item \textbf{Stabilność i skalowalność:} Algorytm wykazuje dużą stabilność, o czym świadczy niewielka różnica między wynikami średnimi a najlepszymi dla każdej z instancji. Potwierdza to, że Local Search jest narzędziem skalowalnym, które w przewidywalny sposób dostosowuje swój czas pracy i jakość rozwiązań do rozmiaru wejściowego problemu optymalizacyjnego.
\end{itemize}

\section*{Zadanie 2}
Zmiana względem zadania 1 polega na przyspieszeniu obliczeń: zamiast sprawdzać wszystkie możliwe sąsiedztwa w jednym kroku, wybieramy najlepsze z n losowych.
Dzięki takiemu podejściu, algorytm jest w stanie szybciej znaleźć względnie dobre rozwiązanie, choć może nie być tak dokładny jak w przypadku pełnego przeszukiwania sąsiedztwa.
% WYNIKI do wykresu:
% Oman avg_res: 140871.20 avg_num_steps: 3010.20 best_res: 136735
% Dżibuti avg res:11354.00 avg_num_steps: 28.80 best_res: 8644
% Irlandia avg_res: 405960.40 avg_num_steps: 12356.60 best_res: 393356
% Egipt avg_res: 296658.20 avg_num_steps: 10888.80best_res: 276573
% Tanzania avg_res: 320460.60 avg_num_steps: 4497.80 best_res: 287197
% Kanada avg_res: 2251333.00 avg_num_steps: 7219.80 best_res: 2196570

\begin{figure}[h!]
    \centering
    % Pierwszy wykres
    \begin{subfigure}[b]{0.45\textwidth}
        \centering
        \includegraphics[width=\textwidth]{/home/wojteq18/Uni-6/Algorytmy_Metaheurystyczne/Lab/Lista_1/charts/21.png}
        \caption{Średnie rozwiązanie}
        \label{fig:wykres1}
    \end{subfigure}
    \hfill % Rozsuwa wykresy na całą szerokość tekstu
    % Drugi wykres
    \begin{subfigure}[b]{0.45\textwidth}
        \centering
        \includegraphics[width=\textwidth]{/home/wojteq18/Uni-6/Algorytmy_Metaheurystyczne/Lab/Lista_1/charts/22.png}
        \caption{Średnia ilość kroków}
        \label{fig:wykres2}
    \end{subfigure}
    \hfill
    % Trzeci wykres
    \begin{subfigure}[b]{0.45\textwidth}
        \centering
        \includegraphics[width=\textwidth]{/home/wojteq18/Uni-6/Algorytmy_Metaheurystyczne/Lab/Lista_1/charts/23.png}
        \caption{Najlepszy rezultat}
        \label{fig:wykres3}
    \end{subfigure}
    
    \caption{Porównanie wyników algorytmu dla wybranych państw}
    \label{fig:wszystkie_wykresy}
\end{figure}

\newpage
\section*{Zadanie 3}
Zmiana względem poprzedniego podejścia polega na wykorzystaniu Minimalnego Drzewa Rozpinającego (MST) jako bazy do budowy trasy. 
Zamiast zaczynać od całkowicie losowej kolejności miast, najpierw wyliczamy MST, a potem na jego podstawie tworzymy cykl startowy. 
Całą procedurę powtarzamy $\lceil \sqrt{n} \rceil$ razy, za każdym razem losując inny wierzchołek, od którego zaczynamy obchodzić drzewo. 
Dzięki temu Local Search nie startuje "w ciemno", tylko dostaje na wejściu już całkiem sensownie poukładaną trasę. 
Pozwala to algorytmowi szybciej dotrzeć do dobrych wyników i sprawdzić różne warianty przejścia tej samej struktury drzewa, 
co zwiększa szansę na znalezienie lepszego rozwiązania końcowego.

% WYNIKI do wykresu:
% Oman avg_res: 93088.56 avg_num_steps: 532.51 best_res: 92155
% Dżibuti avg res:6738.29 avg_num_steps: 12.00 best_res: 6656
% Irlandia avg_res: 225023.76 avg_num_steps: 2161.57 best_res: 224454
% Egipt avg_res: 184201.66 avg_num_steps: 1894.33 best_res: 183905
% Tanzania avg_res: 193849.79 avg_num_steps: 784.04 best_res: 192973
% Kanada avg_res: 1414883.52 avg_num_steps: 1146.57 best_res: 1410394
\begin{figure}[h!]
    \centering
    % Pierwszy wykres
    \begin{subfigure}[b]{0.45\textwidth}
        \centering
        \includegraphics[width=\textwidth]{/home/wojteq18/Uni-6/Algorytmy_Metaheurystyczne/Lab/Lista_1/charts/31.png}
        \caption{Średnie rozwiązanie}
        \label{fig:wykres1}
    \end{subfigure}
    \hfill % Rozsuwa wykresy na całą szerokość tekstu
    % Drugi wykres
    \begin{subfigure}[b]{0.45\textwidth}
        \centering
        \includegraphics[width=\textwidth]{/home/wojteq18/Uni-6/Algorytmy_Metaheurystyczne/Lab/Lista_1/charts/32.png}
        \caption{Średnia ilość kroków}
        \label{fig:wykres2}
    \end{subfigure}
    \hfill
    % Trzeci wykres
    \begin{subfigure}[b]{0.45\textwidth}
        \centering
        \includegraphics[width=\textwidth]{/home/wojteq18/Uni-6/Algorytmy_Metaheurystyczne/Lab/Lista_1/charts/33.png}
        \caption{Najlepszy rezultat}
        \label{fig:wykres3}
    \end{subfigure}
    
    \caption{Porównanie wyników algorytmu dla wybranych państw}
    \label{fig:wszystkie_wykresy}
\end{figure}

\section*{Wnioski}

Na podstawie wyników uzyskanych przy użyciu algorytmu opartego na Minimalnym Drzewie Rozpinającym (MST) sformułowano następujące spostrzeżenia:

\begin{itemize}
    \item \textbf{Naturalna skala problemu:} Podobnie jak w poprzednich zadaniach, wielkość państwa bezpośrednio przekłada się na wartość rozwiązania. Kanada, jako największa instancja, generuje najwyższy koszt trasy, natomiast Dżibuti najniższy, co jest całkowicie zgodne z intuicją.
    
    \item \textbf{Znaczne przyspieszenie zbieżności:} Wykorzystanie MST jako punktu startowego drastycznie zmniejszyło liczbę kroków potrzebnych do znalezienia rozwiązania. 
    
    \item \textbf{Złożoność a liczba kroków:} Mimo że Kanada ma najwyższą wartość rozwiązania, to Irlandia i Egipt wymagają najwięcej kroków poprawy. Potwierdza to wniosek, że trudność optymalizacji zależy od gęstości i układu miast w danym kraju, a nie od samej sumy odległości.
    
    \item \textbf{Wysoka powtarzalność:} Wyniki średnie (\textit{avg\_res}) są bardzo zbliżone do wyników najlepszych (\textit{best\_res}). Dzięki temu, że Local Search startuje z solidnego fundamentu (MST), algorytm jest stabilny i za każdym razem trafia w wysokiej jakości optima lokalne.
\end{itemize}

\end{document} 